\documentclass[11pt]{article}

\usepackage{xparse,environ}
\usepackage{adjustbox}

% \newcommand{\answerDirectory}{solutions}

% Local override, kept so this homework's PDF matches its original exactly:
% the shared macros (src/macros/macros.tex) now define \REPLACEME differently (without the leading space, plus \xspace).
% Defined before ../../macros/macros.tex with \providecommand so a gradescope wrapper's version still wins.
\providecommand{\REPLACEME}{ \textcolor{red}{\ensuremath{\star \cdots \star}}}

../../macros/macros.tex

\inputAnswer{problem_0_answer}

\doHwHeader{Homework \BLUE{3}}

%%%%% EDIT THE FILES NAMED problem_X_answer.tex.  DON'T EDIT THIS FILE

%%%%%%%%%%%%%%%%%%%%%%%%%%%%%%%%%%%%%%%%%%%%%%%%%%%%%%%%%%%% 
%%%%%%%%%%%%%%%%%%%%%%%%%%%%%%%%%%%%%%%%%%%%%%%%%%%%%%%%%%%% 
%%%%%%%%%%%%%%%%%%%%%%%%%%%%%%%%%%%%%%%%%%%%%%%%%%%%%%%%%%%% 
%%%%%%%%%%%%%%%%%%%%%%%%%%%%%%%%%%%%%%%%%%%%%%%%%%%%%%%%%%%% 

\begin{document}


%%%%% EDIT THE FILES NAMED problem_X_answer.tex.  DON'T EDIT THIS FILE 

\paragraph{General instructions.}
% First read \BLUE{Lecture Note 4}, then
Answer the problems given in the rest of this document.
Edit this LaTeX template to add your answers, as described in the file named \texttt{00\_README.txt},
then compile it to produce \BLUE{\texttt{homework\_3.pdf}}\ifGradescope{} for submission to gradescope\fi.

\ifGradescope
For this homework, when you submit to gradescope you don't need to identify the pages
for each problem.  We're using gradescope's ``fixed-length'' grading mode.
To make it work, your answer for each part should fit in the allotted space,
so that each subsequent answer is on the expected page in the expected location.
Your PDF should have
13
pages.  
\fi

\paragraph{Don't know?  Say so.}
If you don't know the answer to a problem or any part of the problem,
we encourage you to simply write \emph{``I don't know''} instead of giving an answer.
(If you write ``I don't know'' you can also add any questions you might have.)

\paragraph{Gentle reminder: homework is for learning.}
As noted in \LN{prelim}, engaging fully with the homeworks is the best way to learn the material.
\emph{The primary role of homeworks is to help you learn the material,
  and the primary role of feedback on homeworks is to help you improve your ideas.}

Please try to explain your ideas as clearly as you can.
The clearer your explanations are, the more useful the feedback on them can be.

\paragraph{Collaboration and citation policy.}
If you are doing this homework as part of a course, follow the course's collaboration and citation policy.
In any case, at the end of each answer, list the people you collaborated with,
and cite every source that your answer uses ideas from
(other than the lecture notes and the sources listed in the lecture notes).
If there are none, write ``none''.

%%%%% EDIT THE FILES NAMED problem_X_answer.tex.  DON'T EDIT THIS FILE 

\vfill
\noindent{\footnotesize\copyrightNotice}

%%% Local Variables:
%%% mode: latex
%%% TeX-master: "homework_2"
%%% End:


%%%%%%%%%%%%%%%%%%%%%%%%%%%%%%%%%%%%%%%%%%%%%%%%%%%%%%%%%%%% 
%%%%%%%%%%%%%%%%%%%%%%%%%%%%%%%%%%%%%%%%%%%%%%%%%%%%%%%%%%%% 
%%%%%%%%%%%%%%%%%%%%%%%%%%%%%%%%%%%%%%%%%%%%%%%%%%%%%%%%%%%% 
%%%%%%%%%%%%%%%%%%%%%%%%%%%%%%%%%%%%%%%%%%%%%%%%%%%%%%%%%%%% 

%%%%% EDIT THE FILES NAMED problem_X_answer.tex.  DON'T EDIT THIS FILE 

%%%%%%%%%%%%%%%% PROBLEM 1  %%%%%%%%%%%%%%%%%

\setcounter{problem}{1}

\newpage 
\problem[1\,(a)] \emph{(Sections 4.1--4.4)}
{\small
  Some recursive algorithm, given any input of size $n$,
  does $\Theta(n\log n)$ work and then recurses on two subproblems of size $n/2$.
  The base case is when $n=1$.
  Suppose the run time $C(n)$ satisfies the following recurrence relation:
  $C(1) = 0$ and $C(n) = 2 C(n/2) + n\log_2 n$ for $n\ge 2$.
  Use the recursion-tree method to find a closed-form $\Theta$ value for $C(n)$,
  then describe your calculations as requested below.
  Assume $n$ is a power of two.
}

\setlength{\ITEMHT}{0.8in}
{
  \injectEnumerate 

  \inputAnswer{problem_1a_answer}
}

\newpage
\problem[1\,(b)]
{\small
  Professor Bo Zo proposes a variant of Karatsuba's algorithm,
  \textsf{BozoMult}, 
  for multiplying large integers.
  Given any two $n$-bit integers $A$ and $B$ (where $n$ is a power of three)
  \textsf{BozoMult}$(A, B)$
  makes eight recursive calls,
  each computing the product of two $n/3$-bit integers.
  In addition to the recursive calls, it takes linear time, $\Theta(n)$
  (just as Karatsuba's algorithm does, outside of its three recursive calls).
  Define $M(n)$ to be the time that  \textsf{BozoMult} takes to multiply two $n$-bit integers.
  State a recurrence relation for $M(n)$,
  then use the recursion-tree method
  to obtain a closed-form $\Theta$ value for $M(n)$.
  Describe your calculations as requested below.
  Assume $n$ is a power of three.
}

\setlength{\ITEMHT}{0.8in}
{
  \injectEnumerate 

  \inputAnswer{problem_1b_answer}
}

\newpage
\problem[1\,(c)]
{\small
  Consider the following ``lopsided'' recurrence relation:
  $T(n) = 1$ for $n\in\{1,2,3,4\}$ and $T(n) = n + T(\lfloor n/4\rfloor) + T(\lfloor 2n/3 \rfloor)$.
  Choose some constant $c$ (as small as you can easily make work)
  then prove by induction that
  $T(n)$ is at most $c n$ for all $n\ge 1$.
  \emph{Hint: Read Section 4.6 through Lemma 4.4.}
}
% \setlength{\ITEMHT}{0.95in}
% {
  % \injectEnumerate 

  \inputAnswer{problem_1c_answer}
% }

%%%%%%%%%%%%%%%%%%%%%%%%%%%%%%%%%%%%%%%%%%%%%%%%%%%%%%%%%%%% 
%%%%%%%%%%%%%%%%%%%%%%%%%%%%%%%%%%%%%%%%%%%%%%%%%%%%%%%%%%%% 
%%%%%%%%%%%%%%%%%%%%%%%%%%%%%%%%%%%%%%%%%%%%%%%%%%%%%%%%%%%% 
%%%%%%%%%%%%%%%%%%%%%%%%%%%%%%%%%%%%%%%%%%%%%%%%%%%%%%%%%%%% 

%%%%% EDIT THE FILES NAMED problem_X_answer.tex.  DON'T EDIT THIS FILE 

%%%%%%%%%%%%%%%% PROBLEM 2  %%%%%%%%%%%%%%%%%

\newpage

{\small
  \problem \emph{(Section 4.7, Exercise 4.10)}
  Read the definition of the \prob{2D Local Maximum} problem in Section 4.7.2.
  Consider the following algorithm (from that section):
  \begin{myframed}
    \begin{steps}
    \item[] {\sf local-max}$(A[1..n, 1..n])$ \comment{We assume $n\ge 1$.}

      \step Let $X$ be the set consisting of the $2n-1$ cells lying in column $(n+1)/2$ or row $(n+1)/2$ of $A$.

      \step Check each cell in $X$ to see if it is a local maximum in $A$.  If so, return such a cell.  Otherwise:

      \step Let $A[i, j]$ be any cell in $X$ with maximum value (among cells in $X$).

      \step Let $A[i', j']$ be a neighbor of $A[i,j]$ in $A$ with $A[i', j'] > A[i,j]$.  (It exists, as $A[i,j]$ is not a local max.)

      \step Partition $A \setminus X$ (the array $A$, minus the cells in $X$) around $X$ into its four
      $(n-1)/2 \times (n-1)/2$ quadrants---the upper left, upper right, lower left, and lower right---in the natural way.

      \step Among those, let $A'$ be the quadrant containing the cell $A[i', j']$.

      \step Recursively compute a local maximum of the subarray $A'$, and return that cell.
    \end{steps}
  \end{myframed}

  Here is an attempted proof of correctness.
  First we prove a utility lemma:
  \begin{lemma}\label{lemma: utility}
    During any execution of the algorithm that reaches Line 7,
    the quadrant $A'$ (as defined in the algorithm Line 6) contains a local maximum of $A$.
  \end{lemma}
  \begin{longFormProof}
    \step Let $i$, $j$, $i'$ and $j'$ be as defined in the algorithm for that execution.  Consider the following process.
    
    \begin{myframed}
      \begin{steps}
      \item[~(i)] Put a token on the cell $A[i', j']$.

      \item[\,(ii)] While the token's cell is not a local maximum in $A$ do:

      \item[(iii)]~~~~Move the token to a neighboring cell with larger value.
      \end{steps}
    \end{myframed}

    \step Each cell in $X$ is no larger than $A[i,j]$ (by Line 4 of the algorithm), and $A[i, j] < A[i',j']$.

    \step So each cell in $X$ is smaller than $A[i', j']$.

    \step So by inspection of the process it never moves the token into $X$.  So the token stays in $A'$.

    \step The process must stop, because each move increases the value of the cell with the token.

    \step So the token ends on a cell in $A'$ that is also a local maximum of $A$.
  \end{longFormProof}

  Here is the rest of the purported proof:

  \begin{lemma}\label{lemma: local max}
    Given any input where $n=2^k-1$ for some integer $k\ge 1$,
    the algorithm gives a correct output.
  \end{lemma}
  \begin{longFormProof}
    \step The proof is by induction on $k$.
    \step Let $P(k)$ be the predicate ``for any input with $n=2^k-1$, the algorithm gives a correct output.''
    \step For any input with $n=1=2^1-1$, the cross $X$ is the entire array, so Line 2 returns a local max.
    \step So $P(1)$ holds.

    \begin{block}[ind hyp]
      \step\label{step: ind} Consider any $k\ge 2$, 
      and any execution of the algorithm on an input $A$ with $n=2^k-1$.
      Assume $P(k-1)$.

      \begin{block}
        \step \underline{Case 1.} Suppose the algorithm returns in Line 2.
        \step By inspection of Line 2, it returns a local max of $A$.
      \end{block}
      \begin{block}
        \step \underline{Case 2.} Otherwise the algorithm returns in Line 7.
        \step The quadrant $A'$ contains a local max of $A$ (by Lemma~\ref{lemma: utility}).
        Note $(n-1)/2 = (2^k-2)/2 = 2^{k-1}-1$.
        \step By Step~\ref{step: ind}, $P(k-1)$ holds, so the recursive call on $A'$ returns a local max of $A'$.
        \step By the previous two steps, the algorithm (Line 7) returns a local max of $A$.
      \end{block}
      \step So, in either case, the algorithm returns a local max of $A$.  That is, it is correct on input $A$.
    \end{block}

    \step By Block~\ref{ind hyp}, for any $k\ge 2$, if $P(k-1)$ holds, then $P(k)$ holds.
    
    \step By Step 4, $P(1)$ holds.

    \step By the previous two steps and induction, $P(k)$ holds for all $k\ge 1$.

    \step That is, the algorithm gives the correct output on any input where $n=2^k-1$ for some integer $k\ge 1$.
  \end{longFormProof}

  \newpage
  {\small
  % \begin{enumerate}[label=(\alph*)]
  % \item
    In the proof of correctness (the two lemmas),
    what is the \emph{first} step in the long-form proof of either lemma
    that asserts a statement that is not necessarily true?
  % \item
    What, precisely, is the mistake in the reasoning in that line?
  % \item
    Give an input with $n=7$ on which the \prob{2D Local Max} algorithm given above fails.
    Describe the execution of the algorithm on that input.
    We encourage you to verify your answer using the Python implementation
    \pycode{divide_and_conquer/local_maximum.py}
    (listed in Section 4.9 of Lecture Note 4).
  % \end{enumerate}
  }
}

\setlength{\ITEMHT}{1.15in}
{
  \injectEnumerate 

  \inputAnswer{problem_2_answer}
}
% \restoreEnumerate

%%%%%%%%%%%%%%%%%%%%%%%%%%%%%%%%%%%%%%%%%%%%%%%%%%%%%%%%%%%% 
%%%%%%%%%%%%%%%%%%%%%%%%%%%%%%%%%%%%%%%%%%%%%%%%%%%%%%%%%%%% 
%%%%%%%%%%%%%%%%%%%%%%%%%%%%%%%%%%%%%%%%%%%%%%%%%%%%%%%%%%%% 
%%%%%%%%%%%%%%%%%%%%%%%%%%%%%%%%%%%%%%%%%%%%%%%%%%%%%%%%%%%% 

%%%%% EDIT THE FILES NAMED problem_X_answer.tex.  DON'T EDIT THIS FILE 

%%%%%%%%%%%%%%%% PROBLEM 3  %%%%%%%%%%%%%%%%%

\newpage
\newcommand{\segment}[2]{#1\!\to\! #2}
\problem \emph{(Exercise 4.11)}
You have been commissioned to design a new bus system that will run along Huntington Avenue.  The bus system must provide service to $N$ stops (numbered $1,2,\ldots,N$) along the eastbound route.  Commuters may begin their trip at any stop $i$ and end at any stop $j$ where $1\le i < j \le N$.  
Here are two naive ways to design the system:
\begin{itemize}
\item You could have a bus run from stop 1 to stop $N$, stopping at every stop in between. The system would be cheap because it would only require $N-1$ route segments, $\{\segment 1 2, \segment 2 3, \ldots, \segment{(N-1)}N\}$, for the entire system. However, a person traveling from stop $1$ to stop $N$ would have to wait while the bus makes $N-1$ stops.

\item You could have, for every pair of stops $i, j$ with $i < j$, a special express bus from $i$ to $j$.
  No commuter would ever have to make more than one stop.
  However, this system would be expensive as it requires ${N \choose 2} = \Theta(N^2)$ route segments,
  $\{\segment i j : 1\le i < j \le N\}$.
  
\end{itemize}

Using divide and conquer, you will find a compromise---a set of at most $N \log_2 N$ route segments, with the property that for any stops $i$ and $j > i$, the \emph{hop length} from $i$ to $j$ in the set is at most 2.
That is, either the (direct) route segment $\segment i j$ is in the set (so the hop length is 1),
or, for some $m$ with $i<m<j$, the set contains the two segments $\segment i m$ and  $\segment m j$
(so the hop length is 2).

Consider any sequence $(\ell, \ell+1, \ldots, u)$ of consecutive stops (where $1\le \ell \le u \le N$).
Let $n = u-\ell+1$ be the number of stops.
We will recursively define a set $S(\ell, u)$ of segments that handles trips from any stop $i$ to any stop $j$
such that $\ell \le i < j \le u$.

For the base case $n=1$, for every stop $\ell$,
define $S(\ell, \ell)$ to be the empty set.
For the base case $n=2$, for every stop $\ell$,
define $S(\ell, \ell+1)$ to be $\{ \segment{\ell}{\ell+1} \}$.
(Then any commuter can get from stop $\ell$ to $\ell+1$ using just the one segment in $S(\ell, \ell+1)$,
so the hop length from $\ell$ to $\ell+1$ in  $S(\ell, \ell+1)$ is 1.)

Now consider any sequence $(\ell, \ell+1, \ldots, u)$ of stops where the number of stops $n=u-\ell+1$ is more than two.
Let $m=\lfloor (\ell+u)/2\rfloor$ be the middle stop.
\begin{enumerate}[label=(\alph*)]
\item 
  Give a definition for a set $Q(\ell, u)$ of segments such that,
  for every pair of stops $(i, j)$ where $i$ is in the first half of the stops ($\ell\le i \le m$)
  and $j$ is in the second half of the stops ($m <j\le u$),
  the hop length from $i$ to $j$ in $Q(\ell, u)$ is at most 2.
  Your set $Q(\ell,  u)$ should contain at most $n-1$ segments (the number of stops minus one).
  This definition should not be recursive.

\item Now give a recursive definition for a set $S(\ell, u)$ of segments such that
  \begin{description}
  \item[(P1)] For \emph{any} pair of stops $(i, j)$ with $\ell\le i < j \le u$,
    the hop length from $i$ to $j$ in $S(\ell, u)$ is at most two.
  \item[(P2)] $S(\ell, u)$ contains at most $n\log_2 n$ segments. 
  \end{description}
  \emph{Hint: Divide the stops into the first and second halves,
    recurse on each half, then add $Q(\ell, u)$.}

\item Let $N(n)$ be the maximum number of segments in any set $S(\ell, u)$ such that $u-\ell+1=n$
  (the set covers $n$ consecutive stops).  State a recurrence relation for $N(n)$.
  Use the recurrence to prove by induction that $N(n) \le n\log_2 n$ for every $n\ge 1$.

\item Prove by induction that, for all $n\ge 1$,
  for any sequence $(\ell,\ell+1, \ldots, u)$ of $n$ consecutive stops (so $u-\ell+1=n$),
  the set $S(\ell, u)$ as you defined it has property (P1) above.
\end{enumerate}

\continued

\setlength{\ITEMHT}{3in}
{
  \injectEnumerate 

  \inputAnswer{problem_3_answer}
}

%%%%%%%%%%%%%%%%%%%%%%%%%%%%%%%%%%%%%%%%%%%%%%%%%%%%%%%%%%%% 
%%%%%%%%%%%%%%%%%%%%%%%%%%%%%%%%%%%%%%%%%%%%%%%%%%%%%%%%%%%% 
%%%%%%%%%%%%%%%%%%%%%%%%%%%%%%%%%%%%%%%%%%%%%%%%%%%%%%%%%%%% 
%%%%%%%%%%%%%%%%%%%%%%%%%%%%%%%%%%%%%%%%%%%%%%%%%%%%%%%%%%%% 

%%%%% EDIT THE FILES NAMED problem_X_answer.tex.  DON'T EDIT THIS FILE 

\end{document}

%%% Local Variables:
%%% mode: latex
%%% TeX-master: t
%%% End:
